\documentclass[12pt]{article}
\usepackage[margin=1in]{geometry}
\usepackage[T1]{fontenc}
\usepackage[utf8]{inputenc}
\usepackage{hyperref}
\usepackage{setspace}

\title{Abstract: CapitolAlpha}
\author{Yinka Vaughan}
\date{April 2026}

\begin{document}
\maketitle
\begin{abstract}
This project examines whether members of the United States Congress and their immediate families, as reported in official legislative financial disclosures, achieve abnormal market returns relative to broad equity benchmarks. Using a data pipeline built in Python, the analysis consolidates Senate and House transactional disclosures from both HTML and PDF sources, harmonizes inconsistent trade records, and enriches the dataset with asset tickers and market-return data. The study focuses on the timing, direction, and magnitude of disclosed trades to identify patterns that may suggest informational advantages or systemic loopholes in current disclosure practices.

Preliminary findings indicate that congressional purchase transactions show stronger positive performance than would be expected by chance over a 90-day holding horizon. These results highlight the importance of transparency and compliance in public-service trading, while also underscoring limitations in the data caused by reporting windows, value bracket ranges, and spouse/trust attribution. By combining quantitative financial analysis with ethical considerations, this work contributes to the broader discussion of legislative accountability and the role of data-driven oversight in democratic institutions.
\end{abstract}

\end{document}
